\documentclass[11pt]{article}
\usepackage{discrete-practice}

\newcommand{\ColloquiumHeader}[3]{%
  \csname Practice@header\endcsname{#1}{Коллоквиум}{#2}{Вопросы и ответы}{}%
  \hypersetup{pdftitle={Дискретная математика: Коллоквиум #2}}%
}
\newcommand{\LectureSource}[2]{\PracticeLink{#1}{#2}}
\newcommand{\QA}[4]{%
  \par\addvspace{.65\baselineskip}\noindent
  \begin{minipage}[t]{\linewidth}
    \raggedright
    \textbf{#1. #2}\par\nobreak
    \vspace{.15\baselineskip}
    #3\par\nobreak
    {\small #4\par}%
  \end{minipage}\par
}
\begin{document}
\ColloquiumHeader{2}{3}
\medskip\noindent\small Под каждым ответом указаны файл лекции и страницы PDF.
Номера весенних лекций соответствуют именам файлов в архиве.
\par\normalsize

\QA{1}{Инцидентность вершин и рёбер}{Вершина и ребро инцидентны, если вершина является одним из концов этого ребра.}{\LectureSource{Весенний архив, лекция 1, стр.~5}{https://drive.google.com/file/d/1M_JKz6itJFag191hicdYI_UGXj5VK4rL/view\#page=5}}

\QA{2}{Смежность вершин и рёбер}{
\begin{enumerate}[label=\arabic*),parsep=0pt]
\item \textbf{Смежные вершины.} Являются концами одного ребра.
\item \textbf{Смежные рёбра.} Имеют общую вершину.
\end{enumerate}
}{\LectureSource{Весенний архив, лекция 1, стр.~5}{https://drive.google.com/file/d/1M_JKz6itJFag191hicdYI_UGXj5VK4rL/view\#page=5}}

\QA{3}{Ребро, ориентированное ребро, петля, кратное ребро}{
\begin{enumerate}[label=\arabic*),parsep=0pt]
\item \textbf{Неориентированное ребро.} Неупорядоченная пара концов $\{u,v\}$.
\item \textbf{Ориентированное ребро (дуга).} Упорядоченная пара $(u,v)$.
\item \textbf{Петля.} Ребро с совпадающими концами.
\item \textbf{Кратные рёбра.} Несколько рёбер, соединяющих одну и ту же пару вершин.
\end{enumerate}
}{\LectureSource{Весенний архив, лекция 1, стр.~5--8}{https://drive.google.com/file/d/1M_JKz6itJFag191hicdYI_UGXj5VK4rL/view\#page=5}}

\QA{4}{Изолированная вершина, висячая вершина и висячее ребро}{
\begin{enumerate}[label=\arabic*),parsep=0pt]
\item \textbf{Изолированная вершина.} Вершина, которой не инцидентно ни одно ребро.
\item \textbf{Висячая вершина.} Вершина, которой инцидентно ровно одно ребро.
\item \textbf{Висячее ребро.} Единственное ребро, инцидентное висячей вершине.
\end{enumerate}
}{\LectureSource{Весенний архив, лекция 1, стр.~8}{https://drive.google.com/file/d/1M_JKz6itJFag191hicdYI_UGXj5VK4rL/view\#page=8}}

\QA{5}{Простой граф, ориентированный граф}{
\begin{enumerate}[label=\arabic*),parsep=0pt]
\item \textbf{Простой граф.} Граф без петель и кратных рёбер.
\item \textbf{Ориентированный граф.} Граф с дугами: порядок концов дуги задаёт её направление.
\end{enumerate}
}{\LectureSource{Весенний архив, лекция 1, стр.~6--9}{https://drive.google.com/file/d/1M_JKz6itJFag191hicdYI_UGXj5VK4rL/view\#page=6}}

\QA{6}{Мультиграф, псевдограф}{В терминологии лекции:
\begin{enumerate}[label=\arabic*),parsep=0pt]
\item \textbf{Мультиграф.} Граф, содержащий кратные рёбра.
\item \textbf{Псевдограф.} Граф, содержащий петли; в нём могут быть и кратные рёбра.
\end{enumerate}
}{\LectureSource{Весенний архив, лекция 1, стр.~9}{https://drive.google.com/file/d/1M_JKz6itJFag191hicdYI_UGXj5VK4rL/view\#page=9}}

\QA{7}{Взвешенный, направленный граф}{
\begin{enumerate}[label=\arabic*),parsep=0pt]
\item \textbf{Взвешенный граф.} Граф, в котором каждому ребру приписан вес.
\item \textbf{Направленный граф.} Орграф без пар противоположных дуг $(u,v)$ и $(v,u)$.
\end{enumerate}
}{\LectureSource{Весенний архив, лекция 1, стр.~6--7}{https://drive.google.com/file/d/1M_JKz6itJFag191hicdYI_UGXj5VK4rL/view\#page=6}}

\QA{8}{Нуль-граф, пустой граф, тривиальный граф}{В терминологии лекции:
\begin{enumerate}[label=\arabic*),parsep=0pt]
\item \textbf{Нуль-граф.} Граф без рёбер.
\item \textbf{Пустой граф.} Граф без вершин.
\item \textbf{Тривиальный граф.} Граф с одной вершиной и без рёбер.
\end{enumerate}
}{\LectureSource{Весенний архив, лекция 1, стр.~7}{https://drive.google.com/file/d/1M_JKz6itJFag191hicdYI_UGXj5VK4rL/view\#page=7}}

\QA{9}{Изоморфизм графов. Помеченный граф}{
\begin{enumerate}[label=\arabic*),parsep=0pt]
\item \textbf{Изоморфизм.} Биекция между множествами вершин, сохраняющая смежность в обе стороны. Равносильно: после переименования вершин графы совпадают и по меткам, и по связям.
\item \textbf{Помеченный граф.} Граф, вершины которого различаются метками.
\end{enumerate}
}{\LectureSource{Весенний архив, лекция 1, стр.~18}{https://drive.google.com/file/d/1M_JKz6itJFag191hicdYI_UGXj5VK4rL/view\#page=18}}

\QA{10}{Способы представления графа}{Диаграмма, списки смежности, матрица смежности \guillemotleft{}вершина на вершину\guillemotright{} и матрица инцидентности \guillemotleft{}вершина на ребро\guillemotright{}.}{\LectureSource{Весенний архив, лекция 1, стр.~13--17}{https://drive.google.com/file/d/1M_JKz6itJFag191hicdYI_UGXj5VK4rL/view\#page=13}}

\QA{11}{Описание кратных рёбер и петель в матрице смежности}{Кратность между различными вершинами записывают в соответствующей ячейке. На диагонали неориентированную петлю учитывают дважды, ориентированную -- один раз.}{\LectureSource{Весенний архив, лекция 1, стр.~14--15}{https://drive.google.com/file/d/1M_JKz6itJFag191hicdYI_UGXj5VK4rL/view\#page=14}}

\QA{12}{Описание кратных рёбер и петель в матрице инцидентности}{Каждому ребру соответствует отдельный столбец, поэтому у параллельных рёбер одинаковые столбцы. В примере лекции неориентированная петля обозначается числом $2$ в строке своей вершины. Для петли орграфа в примере отмечены совпадающие начало и конец ($+1$ и $-1$).}{\LectureSource{Весенний архив, лекция 1, стр.~16--17}{https://drive.google.com/file/d/1M_JKz6itJFag191hicdYI_UGXj5VK4rL/view\#page=16}}

\QA{13}{Степень вершины в неориентированном и ориентированном графах}{
\begin{enumerate}[label=\arabic*),parsep=0pt]
\item \textbf{Неориентированный граф.} Степень -- число концов рёбер, приходящихся на вершину. Петля даёт два конца.
\item \textbf{Орграф, степень входа.} $\deg^-(v)$ -- число дуг, входящих в $v$.
\item \textbf{Орграф, степень исхода.} $\deg^+(v)$ -- число дуг, выходящих из $v$.
\end{enumerate}
}{\LectureSource{Весенний архив, лекция 1, стр.~10--12}{https://drive.google.com/file/d/1M_JKz6itJFag191hicdYI_UGXj5VK4rL/view\#page=10}}

\QA{14}{Лемма о рукопожатиях для неориентированного графа}{$\sum_{v\in V}\deg(v)=2|E|$, потому что у каждого ребра два конца. Отсюда число вершин нечётной степени чётно.}{\LectureSource{Весенний архив, лекция 1, стр.~11}{https://drive.google.com/file/d/1M_JKz6itJFag191hicdYI_UGXj5VK4rL/view\#page=11}}

\QA{15}{Лемма о рукопожатиях для ориентированного графа}{Каждая дуга даёт один вход и один выход. Поэтому $\sum_v\deg^-(v)=\sum_v\deg^+(v)=|E|$, а сумма обеих степеней по всем вершинам равна $2|E|$.}{\LectureSource{Весенний архив, лекция 1, стр.~12}{https://drive.google.com/file/d/1M_JKz6itJFag191hicdYI_UGXj5VK4rL/view\#page=12}}

\QA{16}{Подграф}{$H=(V',E')$ -- подграф $G=(V,E)$, если $V'\subseteq V$, $E'\subseteq E$ и оба конца любого ребра $E'$ входят в $V'$.}{\LectureSource{Весенний архив, лекция 1, стр.~19}{https://drive.google.com/file/d/1M_JKz6itJFag191hicdYI_UGXj5VK4rL/view\#page=19}}

\QA{17}{Надграф}{Граф $H$ является надграфом $G$, если $G$ -- подграф $H$: можно добавить вершины и рёбра, сохранив исходные.}{\LectureSource{Весенний архив, лекция 1, стр.~19}{https://drive.google.com/file/d/1M_JKz6itJFag191hicdYI_UGXj5VK4rL/view\#page=19}}

\QA{18}{Собственный и несобственный подграфы}{
\begin{enumerate}[label=\arabic*),parsep=0pt]
\item \textbf{Собственный подграф.} Непустой подграф, отличающийся от исходного графа.
\item \textbf{Несобственный подграф.} Пустой подграф или сам исходный граф.
\end{enumerate}
}{\LectureSource{Весенний архив, лекция 1, стр.~19}{https://drive.google.com/file/d/1M_JKz6itJFag191hicdYI_UGXj5VK4rL/view\#page=19}}

\QA{19}{Частичный граф}{Подграф, содержащий все вершины исходного графа и некоторое подмножество его рёбер.}{\LectureSource{Весенний архив, лекция 1, стр.~19}{https://drive.google.com/file/d/1M_JKz6itJFag191hicdYI_UGXj5VK4rL/view\#page=19}}

\QA{20}{Регулярный (однородный) граф, полный граф}{
\begin{enumerate}[label=\arabic*),parsep=0pt]
\item \textbf{Регулярный (однородный) граф.} Все вершины имеют одинаковую степень.
\item \textbf{Полный граф.} Простой граф, в котором смежны любые две различные вершины; равносильно, степень каждой вершины равна $n-1$.
\end{enumerate}
}{\LectureSource{Весенний архив, лекция 1, стр.~20}{https://drive.google.com/file/d/1M_JKz6itJFag191hicdYI_UGXj5VK4rL/view\#page=20}}

\QA{21}{Дополнение, объединение, пересечение графов}{
\begin{enumerate}[label=\arabic*),parsep=0pt]
\item \textbf{Дополнение.} Сохраняет вершины простого графа и содержит ровно те рёбра полного графа, которых не было в исходном.
\item \textbf{Объединение.} Объединяет множества вершин и множества рёбер.
\item \textbf{Пересечение.} Пересекает множества вершин и множества рёбер.
\end{enumerate}
}{\LectureSource{Весенний архив, лекция 1, стр.~21}{https://drive.google.com/file/d/1M_JKz6itJFag191hicdYI_UGXj5VK4rL/view\#page=21}}

\QA{22}{Путь. Цепь. Простая цепь}{
\begin{enumerate}[label=\arabic*),parsep=0pt]
\item \textbf{Путь (маршрут).} Последовательность чередующихся вершин и соединяющих их рёбер.
\item \textbf{Цепь.} Маршрут без повторяющихся рёбер.
\item \textbf{Простая цепь.} Цепь, в которой не повторяются и вершины.
\end{enumerate}
}{\LectureSource{Весенний архив, лекция 2, стр.~4}{https://drive.google.com/file/d/1uV-3I-INKuTwh_JXyu06gDkGL9PRQtxT/view\#page=4}}

\QA{23}{Замкнутый путь, цикл и простой цикл}{
\begin{enumerate}[label=\arabic*),parsep=0pt]
\item \textbf{Замкнутый путь.} Начальная и конечная вершины совпадают.
\item \textbf{Цикл.} Замкнутая цепь.
\item \textbf{Простой цикл.} Замкнутая цепь, в которой вершины не повторяются, кроме начала и конца; в определении слайда длина не меньше трёх.
\end{enumerate}
}{\LectureSource{Весенний архив, лекция 2, стр.~4}{https://drive.google.com/file/d/1uV-3I-INKuTwh_JXyu06gDkGL9PRQtxT/view\#page=4}}

\QA{24}{Двудольный граф. Критерий двудольности}{
\begin{enumerate}[label=\arabic*),parsep=0pt]
\item \textbf{Определение.} Вершины можно разбить на две доли без рёбер внутри каждой из них.
\item \textbf{Критерий.} Все простые циклы имеют чётную длину.
\end{enumerate}
}{\LectureSource{Весенний архив, лекция 2, стр.~3}{https://drive.google.com/file/d/1uV-3I-INKuTwh_JXyu06gDkGL9PRQtxT/view\#page=3}}

\QA{25}{Вершинно простой путь, рёберно простой путь}{
\begin{enumerate}[label=\arabic*),parsep=0pt]
\item \textbf{Вершинно простой путь.} Не повторяет вершины; это простая цепь.
\item \textbf{Рёберно простой путь.} Не повторяет рёбра; это цепь.
\end{enumerate}
}{\LectureSource{Весенний архив, лекция 2, стр.~4}{https://drive.google.com/file/d/1uV-3I-INKuTwh_JXyu06gDkGL9PRQtxT/view\#page=4}}

\QA{26}{Лемма о простом пути}{Если есть маршрут из $u$ в $v$, то есть простая цепь между ними. Из маршрута можно последовательно вырезать участки между повторными посещениями одной вершины.}{\LectureSource{Весенний архив, лекция 2, стр.~5}{https://drive.google.com/file/d/1uV-3I-INKuTwh_JXyu06gDkGL9PRQtxT/view\#page=5}}

\QA{27}{Связность в неориентированном графе, компоненты связности. Связный граф}{
\begin{enumerate}[label=\arabic*),parsep=0pt]
\item \textbf{Связность вершин.} Между вершинами существует путь. Это отношение эквивалентности.
\item \textbf{Компоненты связности.} Классы этого отношения.
\item \textbf{Связный граф.} Граф, у которого одна компонента связности.
\end{enumerate}
}{\LectureSource{Весенний архив, лекция 2, стр.~6--7}{https://drive.google.com/file/d/1uV-3I-INKuTwh_JXyu06gDkGL9PRQtxT/view\#page=6}}

\QA{28}{Слабая связность, её компоненты и слабо связный граф}{
\begin{enumerate}[label=\arabic*),parsep=0pt]
\item \textbf{Слабая связность вершин.} Вершины связаны после удаления направлений у всех дуг.
\item \textbf{Компоненты слабой связности.} Компоненты полученного неориентированного графа.
\item \textbf{Слабо связный граф.} Граф, у которого одна компонента слабой связности.
\end{enumerate}
}{\LectureSource{Весенний архив, лекция 2, стр.~8}{https://drive.google.com/file/d/1uV-3I-INKuTwh_JXyu06gDkGL9PRQtxT/view\#page=8}}

\QA{29}{Сильная связность, её компоненты и сильно связный граф}{
\begin{enumerate}[label=\arabic*),parsep=0pt]
\item \textbf{Сильная связность вершин.} Между вершинами есть ориентированные пути в обе стороны.
\item \textbf{Компоненты сильной связности.} Классы этого отношения эквивалентности.
\item \textbf{Сильно связный граф.} Граф, у которого одна компонента сильной связности.
\end{enumerate}
}{\LectureSource{Весенний архив, лекция 2, стр.~8}{https://drive.google.com/file/d/1uV-3I-INKuTwh_JXyu06gDkGL9PRQtxT/view\#page=8}}

\QA{30}{Рёберная двусвязность. Компоненты и минимальная компонента}{
\begin{enumerate}[label=\arabic*),parsep=0pt]
\item \textbf{Рёберная двусвязность.} Между двумя различными вершинами есть два пути без общих рёбер.
\item \textbf{Компонента.} Подграф на классе эквивалентности вершин, с рёбрами внутри этого класса.
\item \textbf{Минимальная компонента.} Одна вершина.
\end{enumerate}
}{\LectureSource{Весенний архив, лекция 2, стр.~10--11}{https://drive.google.com/file/d/1uV-3I-INKuTwh_JXyu06gDkGL9PRQtxT/view\#page=10}}

\QA{31}{Разрезающее множество. Разрез}{
\begin{enumerate}[label=\arabic*),parsep=0pt]
\item \textbf{Разрезающее множество.} Множество рёбер, удаление которого увеличивает число компонент связности.
\item \textbf{Разрез.} Разрезающее множество, минимальное по включению.
\end{enumerate}
}{\LectureSource{Весенний архив, лекция 2, стр.~12}{https://drive.google.com/file/d/1uV-3I-INKuTwh_JXyu06gDkGL9PRQtxT/view\#page=12}}

\QA{32}{Мост: определение и два эквивалентных определения}{
\begin{enumerate}[label=\arabic*),parsep=0pt]
\item Ребро само образует разрез, то есть одноэлементное разрезающее множество.
\item Ребро соединяет разные компоненты рёберной двусвязности.
\item Удаление ребра увеличивает число компонент связности на одну. Если граф был связным, он становится несвязным.
\item В связном графе есть разбиение вершин на две непустые части $U,W$, и любой простой путь из $U$ в $W$ проходит через это ребро.
\item Есть две вершины, соединённые путём, и каждый простой путь между ними проходит через это ребро.
\item Ребро не входит ни в один цикл.
\end{enumerate}
}{\LectureSource{Весенний архив, лекция 2, стр.~12--13}{https://drive.google.com/file/d/1uV-3I-INKuTwh_JXyu06gDkGL9PRQtxT/view\#page=12}}

\QA{33}{Лемма о цикле и мосте}{Ребро является мостом тогда и только тогда, когда оно не лежит ни в одном цикле.}{\LectureSource{Весенний архив, лекция 2, стр.~13}{https://drive.google.com/file/d/1uV-3I-INKuTwh_JXyu06gDkGL9PRQtxT/view\#page=13}}

\QA{34}{Лемма про удаление моста}{Удаление моста увеличивает число компонент связности ровно на одну.}{\LectureSource{Весенний архив, лекция 2, стр.~13}{https://drive.google.com/file/d/1uV-3I-INKuTwh_JXyu06gDkGL9PRQtxT/view\#page=13}}

\QA{35}{Вершинная двусвязность. Компоненты и минимальная компонента}{
\begin{enumerate}[label=\arabic*),parsep=0pt]
\item \textbf{Вершинная двусвязность.} Концы двух рёбер можно соединить двумя путями без общих вершин.
\item \textbf{Компонента (блок).} Подграф из рёбер одного класса эквивалентности и их концов.
\item \textbf{Минимальная компонента.} В примере на слайде -- одно ребро с его концами.
\end{enumerate}
}{\LectureSource{Весенний архив, лекция 2, стр.~14--15}{https://drive.google.com/file/d/1uV-3I-INKuTwh_JXyu06gDkGL9PRQtxT/view\#page=14}}

\QA{36}{Блок: определение через неразделимый подграф и эквивалентное определение}{Два общих описания блока:
\begin{enumerate}[label=\arabic*),parsep=0pt]
\item Граф неразделим, если он связен и не имеет точек сочленения. Блок -- максимальный по включению неразделимый подграф.
\item Компонента вершинной двусвязности: подграф, состоящий из рёбер одного класса вершинной двусвязности и их концов.
\end{enumerate}
Для связного графа хотя бы с тремя вершинами лекция даёт ещё шесть равносильных признаков блока:
\begin{enumerate}[label=\arabic*),parsep=0pt]
\item Любые две вершины лежат на общем цикле.
\item Любая вершина и любое ребро, не являющееся петлёй, лежат на общем цикле.
\item Любые два ребра, не являющиеся петлями, лежат на общем цикле.
\item Для любых двух различных вершин и любого ребра, не являющегося петлёй, найдётся простая цепь между вершинами, проходящая через это ребро.
\item Для любых попарно различных $u,v,w$ есть простая цепь из $u$ в $w$ через $v$.
\item Для любых попарно различных $u,v,w$ есть простая цепь из $u$ в $w$, которая не проходит через $v$.
\end{enumerate}
}{\LectureSource{Весенний архив, лекция 2, стр.~14}{https://drive.google.com/file/d/1uV-3I-INKuTwh_JXyu06gDkGL9PRQtxT/view\#page=14}\enspace\textperiodcentered\enspace \LectureSource{стр.~16}{https://drive.google.com/file/d/1uV-3I-INKuTwh_JXyu06gDkGL9PRQtxT/view\#page=16}\enspace\textperiodcentered\enspace \LectureSource{Весенний архив, лекция 3, стр.~4}{https://drive.google.com/file/d/1hwUOPo2YReiUiJTyAbKLgaB3sNSHQk9c/view\#page=4}}

\QA{37}{Точка сочленения: два эквивалентных определения}{
\begin{enumerate}[label=\arabic*),parsep=0pt]
\item Удаление вершины увеличивает число компонент связности.
\item Вершина принадлежит как минимум двум блокам.
\item В связном графе есть две другие вершины $u,w$, и все простые цепи между ними проходят через данную вершину.
\item В связном графе множество остальных вершин можно разбить на две непустые части так, что любая простая цепь из одной части в другую проходит через данную вершину.
\end{enumerate}
}{\LectureSource{Весенний архив, лекция 3, стр.~2}{https://drive.google.com/file/d/1hwUOPo2YReiUiJTyAbKLgaB3sNSHQk9c/view\#page=2}\enspace\textperiodcentered\enspace \LectureSource{Весенний архив, лекция 3, стр.~3}{https://drive.google.com/file/d/1hwUOPo2YReiUiJTyAbKLgaB3sNSHQk9c/view\#page=3}}

\QA{38}{Лемма о точке сочленения (одна на выбор)}{Леммы со слайда:
\begin{enumerate}[label=\arabic*),parsep=0pt]
\item В связном графе с более чем одной вершиной $v$ является точкой сочленения ровно тогда, когда есть две другие вершины $u,w$, и все простые цепи между ними проходят через $v$. Равносильно: $V\setminus\{v\}$ разбивается на две непустые части, и любой простой путь между частями проходит через $v$.
\item Два различных блока связного графа имеют не более одной общей вершины.
\item В блоке хотя бы с тремя вершинами любые две вершины лежат на общем цикле.
\item В блоке хотя бы с тремя вершинами любая вершина и любое ребро, не являющееся петлёй, лежат на общем цикле.
\end{enumerate}
}{\LectureSource{Весенний архив, лекция 3, стр.~3}{https://drive.google.com/file/d/1hwUOPo2YReiUiJTyAbKLgaB3sNSHQk9c/view\#page=3}}

\QA{39}{Эксцентриситет}{Для вершины $v$ это расстояние до самой далёкой вершины: $\varepsilon(v)=\max_{u\in V}d(v,u)$.}{\LectureSource{Весенний архив, лекция 3, стр.~12}{https://drive.google.com/file/d/1hwUOPo2YReiUiJTyAbKLgaB3sNSHQk9c/view\#page=12}}

\QA{40}{Радиус, диаметр и центр}{
\begin{enumerate}[label=\arabic*),parsep=0pt]
\item \textbf{Радиус.} Минимальный эксцентриситет вершины.
\item \textbf{Диаметр.} Максимальный эксцентриситет, или наибольшее расстояние между двумя вершинами.
\item \textbf{Центр.} Множество вершин, эксцентриситет которых равен радиусу.
\end{enumerate}
}{\LectureSource{Весенний архив, лекция 3, стр.~12}{https://drive.google.com/file/d/1hwUOPo2YReiUiJTyAbKLgaB3sNSHQk9c/view\#page=12}}

\QA{41}{Планарный граф}{
\begin{enumerate}[label=\arabic*),parsep=0pt]
\item Планарный граф можно уложить на плоскости так, чтобы рёбра не пересекались вне общих концов.
\item Плоский граф -- граф с уже выбранной такой укладкой.
\end{enumerate}
}{\LectureSource{Весенний архив, лекция 6, стр.~2}{https://drive.google.com/file/d/1MclkysQbbRJUW55vvEpNKIJ0twWRbLhF/view\#page=2}}

\QA{42}{Грань планарного графа}{Две формулировки для выбранной плоской укладки:
\begin{enumerate}[label=\arabic*),parsep=0pt]
\item Максимальный участок плоскости, любые две точки которого можно соединить кривой, не пересекающей рёбер графа.
\item Одна из областей, на которые граф разбивает плоскость. Неограниченная область называется внешней гранью.
\end{enumerate}
}{\LectureSource{Весенний архив, лекция 6, стр.~2}{https://drive.google.com/file/d/1MclkysQbbRJUW55vvEpNKIJ0twWRbLhF/view\#page=2}}

\QA{43}{Формула Эйлера}{Для связного плоского графа $|V|-|E|+F=2$, где $F$ -- число граней, включая внешнюю.}{\LectureSource{Весенний архив, лекция 6, стр.~3}{https://drive.google.com/file/d/1MclkysQbbRJUW55vvEpNKIJ0twWRbLhF/view\#page=3}}

\QA{44}{Необходимые условия планарности}{Для простого связного планарного графа при $|V|\ge3$: $|E|\le3|V|-6$, откуда следует существование вершины степени не больше пяти. Если граф вдобавок двудольный, $|E|\le2|V|-4$.}{\LectureSource{Весенний архив, лекция 6, стр.~4--6}{https://drive.google.com/file/d/1MclkysQbbRJUW55vvEpNKIJ0twWRbLhF/view\#page=4}}

\QA{45}{Критерий планарности}{
\begin{enumerate}[label=\arabic*),parsep=0pt]
\item \textbf{Критерий Понтрягина -- Куратовского.} Граф планарен тогда и только тогда, когда у него нет подграфов, гомеоморфных $K_5$ или $K_{3,3}$.
\item \textbf{Теорема Вагнера.} Граф планарен тогда и только тогда, когда ни $K_5$, ни $K_{3,3}$ не являются его минорами.
\end{enumerate}
}{\LectureSource{Весенний архив, лекция 6, стр.~7--9}{https://drive.google.com/file/d/1MclkysQbbRJUW55vvEpNKIJ0twWRbLhF/view\#page=7}}

\QA{46}{Дерево (два определения)}{
\begin{enumerate}[label=\arabic*),parsep=0pt]
\item $G$ связен и не имеет циклов.
\item Каждые две различные вершины $G$ соединены ровно одной простой цепью.
\item $G$ связен и $|V|=|E|+1$.
\item $G$ ацикличен и $|V|=|E|+1$.
\item $G$ ацикличен, и при добавлении ребра между любыми двумя несмежными вершинами возникает ровно один простой цикл.
\end{enumerate}
}{\LectureSource{Весенний архив, лекция 3, стр.~17--18}{https://drive.google.com/file/d/1hwUOPo2YReiUiJTyAbKLgaB3sNSHQk9c/view\#page=17}}

\QA{47}{Граф блоков и точек сочленения}{Двудольный граф: в одной доле блоки, в другой точки сочленения. Ребро соединяет блок с каждой входящей в него точкой сочленения. Для связного исходного графа эта конструкция даёт дерево.}{\LectureSource{Весенний архив, лекция 3, стр.~21}{https://drive.google.com/file/d/1hwUOPo2YReiUiJTyAbKLgaB3sNSHQk9c/view\#page=21}}

\QA{48}{Граф компонент рёберной двусвязности}{Каждую компоненту сжимают в вершину; вершины соединяют мостами исходного графа. Для связного исходного графа получается дерево.}{\LectureSource{Весенний архив, лекция 3, стр.~22}{https://drive.google.com/file/d/1hwUOPo2YReiUiJTyAbKLgaB3sNSHQk9c/view\#page=22}}

\QA{49}{Остов графа}{Связный ациклический частичный подграф исходного графа. Он содержит все его вершины.}{\LectureSource{Весенний архив, лекция 4, стр.~8--9}{https://drive.google.com/file/d/1UEcDrRSU2ZYwLPnmCIgZ_105Tf9PJ3hE/view\#page=8}}

\QA{50}{Цикломатическое число}{Число рёбер, которые надо удалить из связного графа, чтобы остался остов. Поскольку у остова $|V|-1$ рёбер, это число равно $|E|-|V|+1$.}{\LectureSource{Весенний архив, лекция 4, стр.~8}{https://drive.google.com/file/d/1UEcDrRSU2ZYwLPnmCIgZ_105Tf9PJ3hE/view\#page=8}\enspace\textperiodcentered\enspace \LectureSource{Весенний архив, лекция 3, стр.~17}{https://drive.google.com/file/d/1hwUOPo2YReiUiJTyAbKLgaB3sNSHQk9c/view\#page=17}}

\QA{51}{Цикл, ассоциированный с остовом}{При добавлении к остову одного ребра вне остова возникает единственный цикл: это ребро вместе с единственным путём в остове между его концами.}{\LectureSource{Весенний архив, лекция 4, стр.~8}{https://drive.google.com/file/d/1UEcDrRSU2ZYwLPnmCIgZ_105Tf9PJ3hE/view\#page=8}\enspace\textperiodcentered\enspace \LectureSource{Весенний архив, лекция 3, стр.~17}{https://drive.google.com/file/d/1hwUOPo2YReiUiJTyAbKLgaB3sNSHQk9c/view\#page=17}}

\QA{52}{Фундаментальная система циклов}{Для выбранного остова -- все циклы, получаемые добавлением к этому остову по отдельности каждого неостовного ребра.}{\LectureSource{Весенний архив, лекция 4, стр.~8}{https://drive.google.com/file/d/1UEcDrRSU2ZYwLPnmCIgZ_105Tf9PJ3hE/view\#page=8}}

\QA{53}{Минимальное остовное дерево}{Остов взвешенного графа с минимальной суммой весов рёбер среди всех остовов.}{\LectureSource{Весенний архив, лекция 4, стр.~11--12}{https://drive.google.com/file/d/1UEcDrRSU2ZYwLPnmCIgZ_105Tf9PJ3hE/view\#page=11}}

\QA{54}{Безопасное ребро}{Ребро вне уже выбранного подграфа $T'$, после добавления которого $T'$ по-прежнему можно дополнить до некоторого минимального остова.}{\LectureSource{Весенний архив, лекция 4, стр.~13}{https://drive.google.com/file/d/1UEcDrRSU2ZYwLPnmCIgZ_105Tf9PJ3hE/view\#page=13}}

\QA{55}{Разрез (разбиение множества вершин)}{Разбиение множества вершин $V$ на две непересекающиеся части $S$ и $T=V\setminus S$.}{\LectureSource{Весенний архив, лекция 4, стр.~14}{https://drive.google.com/file/d/1UEcDrRSU2ZYwLPnmCIgZ_105Tf9PJ3hE/view\#page=14}}

\QA{56}{Ребро, пересекающее разрез}{Ребро, один конец которого находится в $S$, другой -- в $T$.}{\LectureSource{Весенний архив, лекция 4, стр.~14}{https://drive.google.com/file/d/1UEcDrRSU2ZYwLPnmCIgZ_105Tf9PJ3hE/view\#page=14}}

\QA{57}{Лемма о безопасном ребре}{Пусть $T'$ входит в некоторый минимальный остов, и ни одно ребро $T'$ не пересекает выбранный разрез. Любое ребро минимального веса среди пересекающих разрез безопасно для $T'$.}{\LectureSource{Весенний архив, лекция 4, стр.~15}{https://drive.google.com/file/d/1UEcDrRSU2ZYwLPnmCIgZ_105Tf9PJ3hE/view\#page=15}}

\QA{58}{Расстояние между двумя вершинами графа}{Длина кратчайшего маршрута между ними.}{\LectureSource{Весенний архив, лекция 3, стр.~11}{https://drive.google.com/file/d/1hwUOPo2YReiUiJTyAbKLgaB3sNSHQk9c/view\#page=11}}

\QA{59}{Кратчайший путь между двумя вершинами}{Путь между данными вершинами, длина которого равна расстоянию между ними.}{\LectureSource{Весенний архив, лекция 3, стр.~11--12}{https://drive.google.com/file/d/1hwUOPo2YReiUiJTyAbKLgaB3sNSHQk9c/view\#page=11}}

\QA{60}{Эйлеров путь}{Путь, проходящий по каждому ребру графа ровно один раз.}{\LectureSource{Весенний архив, лекция 5, стр.~2}{https://drive.google.com/file/d/17gDjDVg2MeAb0PhQorW7OPuqeSB73gVO/view\#page=2}}

\QA{61}{Отличие эйлерова и полуэйлерова графов}{
\begin{enumerate}[label=\arabic*),parsep=0pt]
\item \textbf{Эйлеров граф.} Содержит эйлеров цикл, то есть замкнутый эйлеров путь.
\item \textbf{Полуэйлеров граф.} Содержит эйлеров путь, но не содержит эйлерова цикла.
\end{enumerate}
}{\LectureSource{Весенний архив, лекция 5, стр.~2}{https://drive.google.com/file/d/17gDjDVg2MeAb0PhQorW7OPuqeSB73gVO/view\#page=2}}

\QA{62}{Эквивалентные определения эйлерова графа}{Для связного графа с более чем одной вершиной эквивалентны:
\begin{enumerate}[label=\arabic*),parsep=0pt]
\item В графе есть эйлеров цикл: замкнутый путь, проходящий по каждому ребру ровно один раз.
\item Степень каждой вершины чётна.
\item Множество рёбер можно разбить на циклы.
\end{enumerate}
Для графа с изолированными вершинами общий критерий приведён в ответе 64.}{\LectureSource{Весенний архив, лекция 5, стр.~4--6}{https://drive.google.com/file/d/17gDjDVg2MeAb0PhQorW7OPuqeSB73gVO/view\#page=4}}

\QA{63}{Теорема о покрытии рёбер графа путями}{При $2N>0$ вершинах нечётной степени в связном графе рёбра можно покрыть $N$ рёберно простыми путями. Связность нужна для применения эйлерова критерия при доказательстве; для произвольного несвязного графа утверждение в таком виде неверно.}{\LectureSource{Весенний архив, лекция 5, стр.~7}{https://drive.google.com/file/d/17gDjDVg2MeAb0PhQorW7OPuqeSB73gVO/view\#page=7}}

\QA{64}{Критерий эйлеровости}{Все вершины имеют чётную степень, а все рёбра принадлежат одной компоненте связности. Изолированные вершины не влияют на возможность обхода рёбер.}{\LectureSource{Весенний архив, лекция 5, стр.~2}{https://drive.google.com/file/d/17gDjDVg2MeAb0PhQorW7OPuqeSB73gVO/view\#page=2}}

\QA{65}{Произвольно вычерчиваемый граф}{Граф произвольно вычерчиваем из $v$, если любую начатую в $v$ цепь можно продолжить до эйлерова цикла.}{\LectureSource{Весенний архив, лекция 5, стр.~8}{https://drive.google.com/file/d/17gDjDVg2MeAb0PhQorW7OPuqeSB73gVO/view\#page=8}}

\QA{66}{Теорема о произвольно вычерчиваемом графе}{Неодноэлементный эйлеров граф произвольно вычерчиваем из $v$ тогда и только тогда, когда каждый его цикл содержит $v$.}{\LectureSource{Весенний архив, лекция 5, стр.~9--10}{https://drive.google.com/file/d/17gDjDVg2MeAb0PhQorW7OPuqeSB73gVO/view\#page=9}}

\QA{67}{Принцип построения произвольно вычерчиваемого графа}{Взять лес и новую вершину $v$. Вершины леса нечётной степени соединить с $v$ нечётным числом рёбер, вершины чётной степени -- чётным числом рёбер. В $v$ можно добавить петли.}{\LectureSource{Весенний архив, лекция 5, стр.~11}{https://drive.google.com/file/d/17gDjDVg2MeAb0PhQorW7OPuqeSB73gVO/view\#page=11}}

\QA{68}{Гамильтонов путь, полугамильтонов граф}{
\begin{enumerate}[label=\arabic*),parsep=0pt]
\item \textbf{Гамильтонов путь.} Проходит через каждую вершину ровно один раз.
\item \textbf{Полугамильтонов граф.} По определению лекции содержит гамильтонов путь.
\end{enumerate}
}{\LectureSource{Весенний архив, лекция 5, стр.~12}{https://drive.google.com/file/d/17gDjDVg2MeAb0PhQorW7OPuqeSB73gVO/view\#page=12}}

\QA{69}{Гамильтонов граф, гамильтонов цикл}{
\begin{enumerate}[label=\arabic*),parsep=0pt]
\item \textbf{Гамильтонов цикл.} Проходит через все вершины; повторяются только совпадающие начало и конец.
\item \textbf{Гамильтонов граф.} Содержит гамильтонов цикл.
\end{enumerate}
}{\LectureSource{Весенний архив, лекция 5, стр.~12}{https://drive.google.com/file/d/17gDjDVg2MeAb0PhQorW7OPuqeSB73gVO/view\#page=12}}

\QA{70}{Теорема Оре}{Если в простом неориентированном графе с $n\ge3$ вершинами для каждой пары несмежных вершин $u,v$ выполнено $\deg(u)+\deg(v)\ge n$, то граф гамильтонов.}{\LectureSource{Весенний архив, лекция 5, стр.~13}{https://drive.google.com/file/d/17gDjDVg2MeAb0PhQorW7OPuqeSB73gVO/view\#page=13}}

\QA{71}{Теорема Дирака}{Если в простом неориентированном графе с $n\ge3$ вершинами каждая вершина имеет степень не меньше $n/2$, то граф гамильтонов.}{\LectureSource{Весенний архив, лекция 5, стр.~13}{https://drive.google.com/file/d/17gDjDVg2MeAb0PhQorW7OPuqeSB73gVO/view\#page=13}}

\QA{72}{Теорема Гуйя-Ури}{Если орграф на $n$ вершинах сильно связен и для каждой вершины выполнены $\deg^+(v)\ge n/2$ и $\deg^-(v)\ge n/2$, то орграф гамильтонов.}{\LectureSource{Весенний архив, лекция 5, стр.~13}{https://drive.google.com/file/d/17gDjDVg2MeAb0PhQorW7OPuqeSB73gVO/view\#page=13}}

\QA{73}{Для любого ли гамильтонова графа выполнены условия теорем о гамильтоновости и почему?}{Нет: условия на слайде достаточные. Из определения следует контрпример: простой цикл $C_5$ гамильтонов, но все степени равны $2$, так что ограничения Дирака и Оре не выполнены. У ориентированного цикла степени входа и выхода равны $1$, поэтому условие из приведённой версии теоремы Гуйя-Ури также не обязательно.}{\LectureSource{Весенний архив, лекция 5, стр.~12--13}{https://drive.google.com/file/d/17gDjDVg2MeAb0PhQorW7OPuqeSB73gVO/view\#page=12}}

\end{document}
